\begin{longtable}{|p{4.3cm}|p{3.2cm}|p{6.5cm}|}
\caption{Phân tách thẩm quyền theo thao tác}
\label{tab:role-permission} \\

\hline
\textbf{Thao tác} & \textbf{Vai trò được phép} & \textbf{Cơ sở của quy định} \\
\hline
\endfirsthead

\hline
\textbf{Thao tác} & \textbf{Vai trò được phép} & \textbf{Cơ sở của quy định} \\
\hline
\endhead

\hline
\endfoot

\hline
\endlastfoot

Cấp và thu hồi quyền quản trị viên
& Chỉ quản trị viên gốc
& Quyền phê duyệt bên phát hành là quyền có hệ quả rộng nhất, nên việc trao quyền đó phải tập trung ở gốc
\\ \hline

Cấp bit ủy quyền cho bên phát hành
& Chỉ quản trị viên gốc
& Bit này cho phép mở rộng chuỗi thẩm quyền; nếu quản trị viên cấp được thì chuỗi có thể mở rộng ngoài kiểm soát của gốc
\\ \hline

Phê duyệt bên phát hành (không kèm bit ủy quyền)
& Quản trị viên gốc, quản trị viên
& Đây là thao tác quản trị thường xuyên, cần phân cấp để vận hành được
\\ \hline

Phê duyệt bên phát hành cấp dưới
& Bên phát hành có bit ủy quyền
& Phạm vi phải là tập con, thời hạn không vượt thời hạn của bên ủy quyền, bit ủy quyền không truyền tiếp
\\ \hline

Đình chỉ, kích hoạt lại, thu hồi bên phát hành
& Quản trị viên gốc, quản trị viên
& Thao tác làm giảm quyền, nên điều kiện tiên quyết lỏng hơn thao tác cấp
\\ \hline

Đề cử người nhận quyền gốc; hủy đề cử
& Chỉ quản trị viên gốc
& Chỉ gốc hiện tại mới có quyền quyết định người kế nhiệm
\\ \hline

Ký nhận quyền gốc
& Địa chỉ đang được đề cử, không chặn theo vai trò
& Tại thời điểm ký nhận, người được đề cử chưa phải gốc; rào chắn thật là điều kiện trùng khớp địa chỉ trong hợp đồng
\\ \hline

Cấp phát chứng thực
& Bên phát hành đang hoạt động, loại nằm trong phạm vi
& Yêu cầu đầy đủ nhất trong hệ thống: cần cả thẩm quyền còn hiệu lực và phạm vi phù hợp
\\ \hline

Thu hồi chứng thực
& Bên đã cấp chứng thực đó
& Chỉ cần quyền sở hữu; không kiểm thẩm quyền còn hiệu lực, theo nguyên tắc bất đối xứng ở Mục 3.4.5
\\ \hline

Đình chỉ chứng thực
& Bên đã cấp chứng thực đó
& Thao tác tạm thời, thuộc quyền của đơn vị cấp
\\ \hline

Kích hoạt lại chứng thực bị đình chỉ
& Chỉ quản trị viên gốc, quản trị viên
& Nếu bên phát hành tự đình chỉ rồi tự kích hoạt lại, thao tác đình chỉ mất giá trị kiểm toán
\\ \hline

Đình chỉ, kích hoạt lại định danh của chính mình
& Chủ sở hữu định danh
& Tự khóa thì tự mở được; hệ thống ghi nhận ai đã đặt trạng thái đình chỉ để phân biệt với trường hợp bị quản trị viên khóa
\\ \hline

Đình chỉ, thu hồi định danh của người khác
& Quản trị viên gốc, quản trị viên
& Định danh bị quản trị viên khóa thì phải quản trị viên mở, tránh việc đối tượng bị xử lý tự giải khóa
\\ \hline

\end{longtable}
